%%%%%%%%%%%%%%%%%%%%%%%%%%
%% Common math notation %%
%%%%%%%%%%%%%%%%%%%%%%%%%%
\newcommand{\Nat}{\mathbb{N}}
\newcommand{\bbrackets}[1]{[\![#1]\!]}
\newcommand{\parto}{\rightharpoonup}

%% Calligraphic letters
\newcommand{\DD}{\mathcal{D}}
\newcommand{\EE}{\mathcal{E}}
\newcommand{\LL}{\mathcal{L}}
\newcommand{\MM}{\mathcal{M}}
\newcommand{\OO}{\mathcal{O}}
\newcommand{\PP}{\mathcal{P}}


%%%%%%%%%%%%%%%%%%%%%%%%%
%% PTS Syntax notation %%
%%%%%%%%%%%%%%%%%%%%%%%%%

%% PTS signature
\newcommand{\St}{\mathcal{S}t}
\newcommand{\Ax}{\mathcal{A}x}
\newcommand{\Ru}{\mathcal{R}u}


%% Names of syntactic categories
\newcommand{\Exp}{\mathsf{Exp}}
\newcommand{\Typ}{\mathsf{Typ}}
\newcommand{\Sub}{\mathsf{Sub}}
\newcommand{\Ctx}{\mathsf{Ctx}}

\newcommand{\Ne}{\mathsf{Ne}}
\newcommand{\Nf}{\mathsf{Nf}}

%% For types
\newcommand{\tppi}[2]{\Pi #1.#2}

%% For terms
\newcommand{\tmcst}[1]{\textbf{#1}}
\newcommand{\tmclo}[2]{[#1]#2}
\newcommand{\tmlam}[1]{\lambda #1}
\newcommand{\tmapp}[2]{#1~#2}

%% For contexts
\newcommand{\ctxempty}{\cdot}
\DeclareDocumentCommand{\ctxext}{m m}{#1, #2}
%% Operations on contexts
\DeclareDocumentCommand{\ctxlookup}{m m m}{#1 \judg^{\scalebox{0.75}{\ensuremath \#}} #2 : #3}


%%For substitutions
\newcommand{\subempty}{\cdot}
\newcommand{\subid}{\texttt{id}}
\newcommand{\subwk}{\texttt{wk}}
\DeclareDocumentCommand{\subext}{m m}{#1, #2}
\newcommand{\subclo}[2]{#1 \circ #2}


%%%%%%%%%%%%%%%
%% Judgments %%
%%%%%%%%%%%%%%%

%% Syntactic judgments
\newcommand{\judg}{\vdash}

\DeclareDocumentCommand{\wfctx}{m}{\judg #1}
\DeclareDocumentCommand{\wfexp}{m m m}{#1 \judg #2 : #3}
\DeclareDocumentCommand{\wfsub}{m m m}{#1 \judg #2 : #3}
\DeclareDocumentCommand{\wftyp}{m m}{#1 \judg #2}

\DeclareDocumentCommand{\eqctx}{m m}{\judg #1 \equiv #2}
\DeclareDocumentCommand{\eqexp}{m m m m}{#1 \judg #2 \equiv #3 : #4}
\DeclareDocumentCommand{\eqsub}{m m m m}{#1 \judg #2 \equiv #3 : #4}
\DeclareDocumentCommand{\eqtyp}{m m m}{#1 \judg #2 \equiv #3}



%% Semantic judgments for completeness
\newcommand{\semjudg}{\vDash}

\DeclareDocumentCommand{\wfctxsem}{m}{\semjudg #1}
\DeclareDocumentCommand{\wfexpsem}{m m m}{#1 \semjudg #2 : #3}
\DeclareDocumentCommand{\wfsubsem}{m m m}{#1 \semjudg #2 : #3}
\DeclareDocumentCommand{\wftypsem}{m m}{#1 \semjudg #2}

\DeclareDocumentCommand{\eqctxsem}{m m}{\semjudg #1 \equiv #2}
\DeclareDocumentCommand{\eqexpsem}{m m m m}{#1 \semjudg #2 \equiv #3 : #4}
\DeclareDocumentCommand{\eqsubsem}{m m m m}{#1 \semjudg #2 \equiv #3 : #4}
\DeclareDocumentCommand{\eqtypsem}{m m m}{#1 \semjudg #2 \equiv #3}



%%%%%%%%%%%%%%%%%%%%%%%%%%%%
%% PTS Semantics notation %%
%%%%%%%%%%%%%%%%%%%%%%%%%%%%

%% Names of semantic categories
\newcommand{\Dom}{\mathsf{D}}
\newcommand{\Dne}{\Dom^{\Ne}}
\newcommand{\Dnf}{\Dom^{\Nf}}
\newcommand{\Env}{\mathsf{Env}}

%% Neutral domain expressions
\DeclareDocumentCommand{\dapp}{m m}{#1~#2}

%% Domain expressions
\DeclareDocumentCommand{\dcst}{m}{\textbf{#1}}
\DeclareDocumentCommand{\dreflect}{m m}{\Uparrow^{#1}\!#2}
\DeclareDocumentCommand{\dpi}{m m m}{\Pi#1.#2(#3)}
\DeclareDocumentCommand{\dlam}{m m}{\lambda #1(#2)}

%% Normal domain expressions
\DeclareDocumentCommand{\dreify}{m m}{\Downarrow^{#1}\!#2}

%% Environments
\DeclareDocumentCommand{\denvempty}{}{\cdot}
\DeclareDocumentCommand{\denvext}{m m}{#1, #2}


%% Evaluation
\DeclareDocumentCommand{\evalexp}{m m m}{{\bbrackets{#1}}^{\Exp}(#2) \searrow #3}
\DeclareDocumentCommand{\evalsub}{m m m}{{\bbrackets{#1}}^{\Sub}(#2) \searrow #3}
\DeclareDocumentCommand{\evalapp}{m m m}{#1 \cdot #2 \searrow #3}
%% Functional versions of evaluation judgments
\DeclareDocumentCommand{\evalexpfunc}{m m}{{\bbrackets{#1}}^{\Exp}(#2)}
\DeclareDocumentCommand{\evalsubfunc}{m m}{{\bbrackets{#1}}^{\Sub}(#2)}
\DeclareDocumentCommand{\evalappfunc}{m m}{#1 \cdot #2}


%% Readbacks
\DeclareDocumentCommand{\rbne}{O{i} m m}{\mathsf{R}_{#1}^{\Ne}(#2) \searrow #3}
\DeclareDocumentCommand{\rbnf}{O{i} m m}{\mathsf{R}_{#1}^{\Nf}(#2) \searrow #3}
\DeclareDocumentCommand{\rbtyp}{O{i} m m}{\mathsf{R}_{#1}^{\Typ}(#2) \searrow #3}
\DeclareDocumentCommand{\rbenv}{O{i} m m}{\mathsf{R}_{#1}^{\Env}(#2) \searrow #3}
%% Functional readbacks
\DeclareDocumentCommand{\rbnefunc}{O{i} m}{\mathsf{R}_{#1}^{\Ne}(#2)}
\DeclareDocumentCommand{\rbnffunc}{O{i} m}{\mathsf{R}_{#1}^{\Nf}(#2)}
\DeclareDocumentCommand{\rbtypfunc}{O{i} m}{\mathsf{R}_{#1}^{\Typ}(#2)}
\DeclareDocumentCommand{\rbenvfunc}{O{i} m}{\mathsf{R}_{#1}^{\Env}(#2)}




%% Operations on environments
\DeclareDocumentCommand{\denvlookup}{m m}{#1(#2)}
\DeclareDocumentCommand{\denvconcat}{m m}{#1;#2}
\DeclareDocumentCommand{\denvsize}{m}{|#1|}

\DeclareDocumentCommand{\denvinit}{m}{\rho^*_{#1}}

\DeclareDocumentCommand{\nbeexp}{m m m m}{\texttt{nbe}^{\Exp}(#1 \judg #2 : #3) \searrow #4}
\DeclareDocumentCommand{\nbeexpfunc}{m m m}{\texttt{nbe}^{\Exp}(#1 \judg #2 : #3)}



%%% PER Model %%%

\newcommand{\per}[1]{\mathcal{#1}}

\newcommand{\perne}{Ne}
\newcommand{\pernf}{Nf}
\newcommand{\pertyp}{Typ}

\newcommand{\perneu}{Neu}

\DeclareDocumentCommand{\persort}{O{\tmcst s}}{\mathcal{S}ort_{#1}}

\newcommand{\perctx}{\mathcal{C}tx}


\DeclareDocumentCommand{\perrel}{m m m}{#1 = #2 \in #3}


%% Candidate spaces
\DeclareDocumentCommand{\perleast}{m}{\underline{#1}}
\DeclareDocumentCommand{\permost}{m}{\overline{#1}}

\DeclareDocumentCommand{\realizes}{}{\vdash\!\!\vdash}



%%%%%%%%%%%%%%%%%%%%%%%%
%% Names of languages %%
%%%%%%%%%%%%%%%%%%%%%%%%
\newcommand{\cocon}{\textsc{Cocon}\xspace}
\newcommand{\stlc}{\textsc{STLC}\xspace}
\newcommand{\lf}{\textsc{LF}\xspace}
\newcommand{\mltt}{\textsc{MLTT}\xspace}
\newcommand{\coc}{\textsc{CoC}\xspace}
\newcommand{\sysf}{\textsc{System F}\xspace}
\newcommand{\cycf}{\textsc{Cyclic F}\xspace}
\newcommand{\agda}{\textsc{Agda}\xspace}
\newcommand{\rocq}{\textsc{Rocq}\xspace}
\newcommand{\lean}{\textsc{Lean}\xspace}



%%%%%%%%%%%%%
%% Editing %%
%%%%%%%%%%%%%
\newcommand{\AG}[1]{{\color{green!60!black}AG: #1}}
\newcommand{\BP}[1]{{\color{red}BP: #1}}

\newcommand{\mhighlight}[1]{\colorbox{light-gray}{\ensuremath{#1}}}
